\subsection{Stone Pile за $\mathcal{O}(N\cdot 2^N)$}

\subsubsection{Описание решения}

В задаче используется ограничение \( 0\leq N\leq 20 \), что намекает на использование алгоритма высокой сложности. Следовательно, можно попробовать перебрать всевозможные варианты неупорядоченных распределений элементов между двумя группами --- их не более \( 2^{20}\approx 10^6 \).

Для оптимизации использования памяти можно воспользоваться битовыми операциями. В операционных системах чаще всего выделяется \( 32 \) бит для \textit{int}, поэтому он подойдёт для ограничений данной задачи.

Таким образом, на хранение вариантов распределения будет выделено хотя бы \( 2^{20}\cdot 2^5\cdot 2^{-23}=4 \) МиБ, что допустимо для ограничений данной задачи.

Далее для каждого сгенерированного варианта при помощи битовых сдвигов вычисляем сумму элементов каждой группы. Если \( i \)-ый бит справа равен единице, то элемент принадлежит первой группе, иначе --- второй.

После мы вычисляем модуль разности сумм и обновляем глобальный минимум в случае необходимости.

\subsubsection{Асимптотическая сложность}

\textit{Временная сложность решения} --- \( \mathcal{O}(N\cdot 2^N) \). После считывания массива генерируются всевозможные варианты неупорядоченного распределения элементов по двум группам. Затем, для каждого варианта вычисляется искомая разность сумм двух групп за линейное время и обновляется глобальный минимум.

\textit{Пространственная сложность решения} --- \( \mathcal{O}(2^N) \). Всевозможные варианты распределения камней между двумя кучами хранятся в массиве длины \( 2^N \), исходные данные хранятся в массиве длины \( N \). 

\subsubsection{Листинг кода}

\begin{lstlisting}
#include <iostream>
#include <vector>

std::vector<int> add_options(std::vector<int> options) {
  std::vector<int> new_options;
  for (int o : options) {
    new_options.push_back(o << 1 | 1);
    new_options.push_back(o << 1);
  }
  return new_options;
}

int main() {
  int N;
  std::cin >> N;
  std::vector<int> arr;
  for (int i = 0; i < N; i++) {
    int item;
    std::cin >> item;
    arr.push_back(item);
  }
  std::vector<int> opts = {0, 1};
  for (int i = 0; i < N - 1; i++) {
    opts = add_options(opts);
  }
  long ans = 2000000;
  for (int o : opts) {
    long sum1 = 0, sum2 = 0, curr;
    for (int i = 0; i < N; i++) {
      if (o >> i & 1) {
        sum1 += arr[i];
      } else {
        sum2 += arr[i];
      }
    }
    curr = abs(sum2 - sum1);
    if (curr < ans) {
      ans = curr;
    }
  }
  std::cout << ans;
  return 0;
}
\end{lstlisting}